\documentclass[12pt]{article}
\usepackage[T1]{fontenc}
\usepackage[utf8]{inputenc}
\usepackage{lmodern}
\usepackage{textcomp}
\usepackage{enumitem}
\usepackage{geometry}
\geometry{margin=1in}
\begin{document}
\begin{center}
Answer Key - Mathematics - K-12 Level Assessment 22 \
\small (Answers are ordered exactly as in the question paper.)
\end{center}

\section*{Multiple Choice}
\begin{enumerate}[label=\arabic*.]
\item \textbf{Answer:} B
\\ \textit{Options:} A) 1 meter; B) 25 meters; C) 1 kilometer; D) 25 kilometers
\\ \textit{Note:} The most reasonable estimate of the length of a city's swimming pool is 25 meters because standard competitive swimming pools are typically 25 meters in length, making this a common size for municipal facilities.
\item \textbf{Answer:} D
\\ \textit{Options:} A) 10; B) 45; C) 125; D) 25
\\ \textit{Note:} The area of a right triangle can be calculated using the formula \( \text{Area} = \frac{1}{2} \times \text{base} \times \text{height} \), and since the triangle is isosceles with a 45-degree angle, the base and height both equal \( 5 \) inches (half of the hypotenuse), resulting in an area of \( \frac{1}{2} \times 5 \times 5 = 25 \) square inches.
\item \textbf{Answer:} B
\\ \textit{Options:} A) -4; B) 0; C) 2; D) 4
\\ \textit{Note:} The correct answer is 0 because, when simplified, the equation 3 x - 4(x - 2) + 6 x - 8 = 0 reduces to 0 = 0, which indicates that x must equal 0 for the equation to hold true.
\item \textbf{Answer:} D
\\ \textit{Options:} A) 12; B) 3; C) 5; D) 6
\\ \textit{Note:} The correct answer is 6 because selecting 5 students from a group of 6 is equivalent to choosing 1 student to leave out, which can be done in 6 different ways.
\item \textbf{Answer:} C
\\ \textit{Options:} A) (-4,4); B) [0,4]; C) (-inf, -1) U (4, inf); D) (-inf, -1) U (-1, 4) U (4, inf)
\\ \textit{Note:} The domain of the function \( f(x) = \frac{x+6}{\sqrt{x^2-3 x-4}} \) is restricted to values of \( x \) that make the denominator positive, which occurs for \( x < -1 \) and \( x > 4 \), thus the domain is \( (-\infty, -1) \cup (4, \infty) \).
\end{enumerate}

\section*{True / False}
\begin{enumerate}[label=\arabic*.]
\setcounter{enumi}{5}
\item \textbf{Answer:} False
\\ \textit{Options:} A) True; B) False
\\ \textit{Note:} The greatest common factor of 42 and 84 is 42, as it is the largest number that divides both evenly.
\item \textbf{Answer:} False
\\ \textit{Options:} A) True; B) False
\\ \textit{Note:} The statement is false because the smallest positive integer n such that 1/n is a terminating decimal and contains the digit 9 is actually 90, not 256.
\item \textbf{Answer:} False
\\ \textit{Options:} A) True; B) False
\\ \textit{Note:} The value of y that satisfies the equation y + 2.9 = 11 is actually 8.1, not 9.1, because when you subtract 2.9 from 11, you get 8.1.
\item \textbf{Answer:} True
\\ \textit{Options:} A) True; B) False
\\ \textit{Note:} The value of N is 36 because it is the solution to the equation when all terms are correctly evaluated and balanced.
\item \textbf{Answer:} True
\\ \textit{Options:} A) True; B) False
\\ \textit{Note:} There are 28 days in February (in a common year), and since there are 7 days in a week, 28 divided by 7 equals 4, making the statement true.
\end{enumerate}

\section*{Long Form Response}
\begin{enumerate}[label=\arabic*.]
\setcounter{enumi}{10}
\item \textbf{Answer:} For $(ax + b)(x^5 + 1) - (5 x + 1)$ to be divisible by $x^2 + 1,$ it must be equal to 0 at the roots of $x^2 + 1 = 0,$ which are $\pm i.$ For $x = i,$ \begin{align*} (ax + b)($x^5$ + 1) - (5 x + 1) \&= (ai + b)(i + 1) - (5 i + 1) \\ \&= -a + ai + bi + b - 5 i - 1 \\ \&= (-a + b - 1) + (a + b - 5)i. \end{align*}Then we must have $-a + b - 1 = a + b - 5 = 0.$ Solving, we find $(a,b) = \boxed{(2,3)}.$ For these values, \[(ax + b)(x^5 + 1) - (5 x + 1) = 2 x^6 + 3 x^5 - 3 x + 2 = (x^2 + 1)(2 x^4 + 3 x^3 - 2 x^2 - 3 x + 2).\]
\\ \textit{Note:} correct because for the polynomial to be divisible by \(x^2 + 1\), it must evaluate to zero at the roots \(x = i\) and \(x = -i\), leading to a system of equations that determines the values of \(a\) and \(b\) as \(2\) and \(3\) respectively.
\item \textbf{Answer:} Jimmy and Kima are planning a three-day road trip with specific mileages of 200 miles, 450 miles, and 200 miles for each day. To estimate the total distance they will travel, we can simply add these mileages together: 200 + 450 + 200. This calculation gives us a total distance of 850 miles. This estimation method is straightforward because it involves basic addition, which is a fundamental skill in mathematics, allowing us to easily determine the overall distance for the trip.
\\ \textit{Note:} correct because it accurately sums the daily mileages of 200, 450, and 200 miles using basic addition, resulting in a total distance of 850 miles for the three-day road trip.
\end{enumerate}

\vfill
\noindent\textit{End of Answer Key}
\end{document}